\documentclass[11pt]{article}
\usepackage[margin=1in]{geometry}
\usepackage[T1]{fontenc}
\usepackage{lmodern}
\usepackage{microtype}
\usepackage{amsmath,amssymb,amsthm,mathtools}
\usepackage{booktabs,array}
\usepackage{enumitem}
\usepackage{float}
\usepackage{xcolor}
\usepackage[colorlinks=true,linkcolor=blue!55!black,citecolor=blue!55!black,urlcolor=blue!55!black]{hyperref}

\newtheorem{theorem}{Theorem}[section]
\newtheorem{proposition}[theorem]{Proposition}
\newtheorem{lemma}[theorem]{Lemma}
\newtheorem{corollary}[theorem]{Corollary}
\newtheorem{remark}[theorem]{Remark}
\newcommand{\HS}{\mathrm{HS}}
\newcommand{\rank}{\operatorname{rank}}
\newcommand{\one}{\mathbf 1}

\hypersetup{
 pdftitle={An Exact Rank--Norm Phase Diagram for Three-Level Self-Commutator Targets},
 pdfauthor={Lluis Eriksson},
 pdfsubject={A complete Horn phase diagram for minimum-norm inverse self-commutators in a dimension-eight spectral cone},
 pdfkeywords={self-commutator, Horn inequalities, inverse commutator, optimal rank, phase diagram, inertia}
}

\title{An Exact Rank--Norm Phase Diagram for\\
Three-Level Self-Commutator Targets}
\author{Lluis Eriksson}
\date{24 August 2026}

\begin{document}
\maketitle

\begin{abstract}
For a traceless Hermitian matrix $F$, let
\[
 \kappa_d(F)=\frac12\min\{\|C\|_{\HS}^2:CC^*-C^*C=2F\}.
\]
We determine this inverse self-commutator cost, the cost under rank at most
four, and the minimum rank on the optimal face for the entire dimension-eight
spectral cone
\[
 \lambda(b,c)=(4c-3b,b,b,b,-c,-c,-c,-c),\qquad 0\le b\le c.
\]
The unrestricted value has two chambers,
\[
 \kappa_8=\begin{cases}10c-7b,&0\le b\le c/2,\\
                         9c-5b,&c/2\le b\le c,
          \end{cases}
\]
whereas the rank-four value is $10c-6b$ throughout.  Hence the exact rank
penalty is the tent function $\min\{b,c-b\}$.  Every interior target has
inertia $(4,4)$ but requires rank five at minimum norm; rank four returns only
on the boundary rays.  The relative rank-four tax is at most $14/13$, with
 equality exactly at $b=c/2$.  Sparse Horn duals certify all three affine
pieces.  We also give a genuinely two-valued dimension-nine target with
spectrum $(5,5,5,5,-4,-4,-4,-4,-4)$: its unrestricted cost is $29$ and first
occurs at rank six, whereas the rank-five cost is $30$.  Thus three spectral
levels are not necessary for the rank gap.  Dependency-free exact replays
verify all $8{,}759$ dimension-eight and $39{,}724$ dimension-nine constraints.
This gives a complete positive-dimensional phase diagram, not a continuity
consequence from one example.  We do not claim that dimension eight is the
smallest dimension where incompatibility occurs.
\end{abstract}

\section{Introduction}

Every finite-dimensional trace-zero matrix is a commutator
\cite{albert1957,thompson1958}.  For Hermitian targets, representing the
matrix as one self-commutator imposes a more rigid inverse problem: among all
$C$ with $CC^*-C^*C=2F$, how small can $\|C\|_{\HS}$ be, and what rank is
needed to attain that minimum?  This is distinct from forward Frobenius
commutator bounds \cite{bottcher2005,bottcher2008} and from unrestricted or
infinite-dimensional commutator questions \cite{weiss1980,fan1980}.

There is an elementary tension.  Inertia forces every implementing factor to
have rank at least $\max\{n_+(F),n_-(F)\}$, while the norm optimum is governed
by the Horn polytope for a difference of two isospectral positive matrices.
The two demands agree in exact solutions through dimensions three, four and
five \cite{erikssonQutrit,erikssonFour,erikssonFive}, and also on certain
dimension-free one-spike strata.  They need not agree in general: an isolated
dimension-eight calculation first revealed a rank-five optimum above a
balanced inertia bound of four.  An isolated point, however, does not reveal
the chamber geometry, the transition mechanism, or the sharp size of the
penalty.

This paper replaces that point calculation by a complete cone theorem.  The
family $\lambda(b,c)$ is two-dimensional before scaling and one-dimensional
projectively.  The proof is finite and transparent: explicit
squared-singular-value spectra give the upper bounds, sparse dual combinations
give the lower bounds, and exact endpoint verification proves every Horn
inequality over each chamber.  As a separate robustness theorem, a two-valued
target in dimension nine has unrestricted optimum rank six although inertia
forces only rank five.

We do not claim Horn feasibility, the inertia obstruction, self-commutator
existence, the earlier compatible cases, or dimensional minimality as new.
The contribution is the exhaustive parametric phase diagram, its exact
rank--norm tradeoff, and the two-valued separation.  The phase has three
notable features:
\begin{enumerate}[leftmargin=2em]
 \item two unrestricted Horn chambers separated at $b=c/2$;
 \item a single affine rank-four branch, separated from the optimum at every
 interior point by $\min\{b,c-b\}$; and
 \item a sharp relative tax $14/13$ at the chamber wall.
\end{enumerate}

\section{Inverse cost and Horn reduction}

Let $F=F^*\in M_d(\mathbb C)$ be traceless and set
\begin{equation}
 \kappa_d(F)=\frac12\min\{\|C\|_{\HS}^2:CC^*-C^*C=2F\}.
 \label{eq:kappa}
\end{equation}
Finite-dimensional self-commutator existence follows, for example, from a
weighted shift after a zero-sum ordering of the eigenvalues
\cite{fan1980,beltita2014}.  Write
\[
 r_0(F)=\max\{n_+(F),n_-(F)\},\qquad
 r_*(F)=\min\{\rank C:C\text{ attains \eqref{eq:kappa}}\}.
\]

\begin{lemma}[inertia obstruction]\label{lem:inertia}
Every feasible $C$ satisfies $\rank C\ge r_0(F)$.
\end{lemma}

\begin{proof}
Put $P=CC^*$ and $Q=C^*C$.  The quadratic form of $P-Q=2F$ is
nonpositive on $\ker P$, so its positive inertia is at most $\rank P$.
Likewise its negative inertia is at most $\rank Q$.  Both ranks equal
$\rank C$.
\end{proof}

We recall the spectral linear program used below.

\begin{proposition}[Horn formulation]\label{prop:horn}
Let $\lambda_1\ge\cdots\ge\lambda_d$ be the eigenvalues of $F$.  For
$p_1\ge\cdots\ge p_{d-1}\ge0$, set
\[
 \alpha(p)=(p_1,\ldots,p_{d-1},0),\quad
 \beta(p)=(0,-p_{d-1},\ldots,-p_1),\quad \gamma=2\lambda.
\]
Then
\begin{equation}
 \kappa_d(F)=\frac12\min\left\{\sum_{j=1}^{d-1}p_j:
 (\alpha(p),\beta(p),\gamma)\text{ satisfies every Horn inequality}\right\}.
 \label{eq:hornlp}
\end{equation}
The least cost among factors of rank at most $r$ is obtained by adding
$p_{r+1}=\cdots=p_{d-1}=0$.
\end{proposition}

\begin{proof}
For a feasible factor, $P=CC^*$ and $Q=C^*C$ are positive semidefinite,
isospectral, and $P+(-Q)=2F$.  At a minimum their smallest common eigenvalue
is zero: if it were $s>0$, the polar decomposition would permit subtracting
$s\one$ from the common positive spectrum while preserving $P-Q$ and reducing
the trace.  Horn's theorem \cite{horn1962,klyachko1998,knutson1999,fulton2000}
then yields \eqref{eq:hornlp}.

Conversely, Horn feasibility supplies positive $P,Q$ with the displayed
isospectral data and $P-Q=2F$.  If $Q=U^*PU$, then $C=P^{1/2}U$ satisfies
$CC^*=P$ and $C^*C=Q$.  Its rank is the number of positive $p_j$.
\end{proof}

For $(I,J,K)\in T_r^8$ we use the convention
\begin{equation}
 \sum_{i\in I}\alpha_i+\sum_{j\in J}\beta_j
 \ge \sum_{k\in K}\gamma_k.
 \label{eq:horn}
\end{equation}
The recursive $T$-sets are those in Fulton's exposition \cite{fulton2000}.

\section{The complete three-level phase diagram}

For $0\le b\le c$, define $F_{b,c}$ to be any Hermitian matrix with ordered
trace-zero spectrum
\begin{equation}
 \lambda(b,c)=(4c-3b,b,b,b,-c,-c,-c,-c).
 \label{eq:family}
\end{equation}
The case $c=0$ is the zero target, so statements about rank below assume
$c>0$.  The interior $0<b<c$ has inertia $(4,4)$.

\begin{theorem}[exact rank--norm phase diagram]\label{thm:phase}
For every $c>0$ and $0\le b\le c$,
\begin{equation}
 \kappa_8(F_{b,c})=
 \begin{cases}
  10c-7b,&0\le b\le c/2,\\
  9c-5b,&c/2\le b\le c.
 \end{cases}
 \label{eq:globalvalue}
\end{equation}
The minimum cost under $\rank C\le4$ is
\begin{equation}
 \kappa_8^{(\le4)}(F_{b,c})=10c-6b.
 \label{eq:facevalue}
\end{equation}
Consequently
\begin{equation}
 \kappa_8^{(\le4)}-\kappa_8=\min\{b,c-b\}.
 \label{eq:tent}
\end{equation}
For $0<b<c$, the minimum rank on the global optimal face is exactly five.
At $b=0$ and $b=c$, a rank-four minimizer exists.
\end{theorem}

The three primal branches are
\begin{align}
 p^{L}&=(8c-6b,6c-6b,4c-4b,2c,2b,0,0),
 &&0\le b\le c/2, \label{eq:plow}\\
 p^{H}&=(8c-6b,4c-2b,2c,2c,2c-2b,0,0),
 &&c/2\le b\le c, \label{eq:phigh}\\
 p^{4}&=(8c-6b,6c-4b,4c-2b,2c,0,0,0),
 &&0\le b\le c. \label{eq:pface}
\end{align}
Their half-traces are respectively $10c-7b$, $9c-5b$, and $10c-6b$.

\subsection{Sparse lower certificates}

The following tables list only the specialized inequalities needed for the
dual proofs.  Their Horn triples are checked recursively by the exact replay.
Write $[8]=\{1,\ldots,8\}$.

\begin{table}[H]
\centering
\scriptsize
\caption{Global dual on $0\le b\le c/2$.  Multipliers are shown explicitly.}
\label{tab:lowdual}
\begin{tabular}{@{}c p{6.2cm} c c@{}}
\toprule
weight & $(I;J;K)$ or elementary row & inequality & RHS\\
\midrule
$1/2$ & $(\{1\};\{1\};\{1\})$ & $p_1\ge8c-6b$ & $8c-6b$\\
$1/2$ & $(\{1,2\};\{1,8\};\{1,8\})$ & $p_2\ge6c-6b$ & $6c-6b$\\
$1/2$ & $(\{1,2,3,5,6,7\};\{1,2,3,5,7,8\};\{1,3,4,6,7,8\})$
& $p_3-p_4+p_5\ge2c-2b$ & $2c-2b$\\
$1$ & $([8]\!\setminus\!\{8\};[8]\!\setminus\!\{5\};[8]\!\setminus\!\{5\})$
& $p_4\ge2c$ & $2c$\\
$1/2,1$ & $p_6-p_7\ge0$, $p_7\ge0$ & positivity completion & $0$\\
\bottomrule
\end{tabular}
\end{table}

\begin{table}[H]
\centering
\small
\caption{Global dual on $c/2\le b\le c$.  The multipliers are
$1/2,1/2,1/2,1,1$.}
\label{tab:highdual}
\begin{tabular}{@{}c p{6.1cm} c@{}}
\toprule
$r$ & $(I;J;K)$ or elementary row & specialized inequality\\
\midrule
1 & $(\{1\};\{1\};\{1\})$ & $p_1\ge8c-6b$\\
4 & $(\{1,2,5,6\};\{1,2,5,8\};\{1,4,7,8\})$
& $p_2-p_4+p_5+p_6-p_7\ge4c-4b$\\
7 & $([8]\!\setminus\!\{8\};[8]\!\setminus\!\{6\};[8]\!\setminus\!\{6\})$
& $p_3\ge2c$\\
7 & $([8]\!\setminus\!\{8\};[8]\!\setminus\!\{5\};[8]\!\setminus\!\{5\})$
& $p_4\ge2c$\\
-- & $p_7\ge0$ & positivity\\
\bottomrule
\end{tabular}
\end{table}

\begin{table}[H]
\centering
\small
\caption{Rank-four dual on the face $p_5=p_6=p_7=0$.  Every multiplier is
$1/2$.}
\label{tab:facedual}
\begin{tabular}{@{}c p{6.1cm} c@{}}
\toprule
$r$ & $(I;J;K)$ & specialized face inequality\\
\midrule
1 & $(\{1\};\{1\};\{1\})$ & $p_1\ge8c-6b$\\
3 & $(\{1,2,5\};\{1,2,8\};\{1,4,8\})$ & $p_2\ge6c-4b$\\
5 & $(\{1,2,3,5,6\};\{1,2,3,7,8\};\{1,3,4,7,8\})$ & $p_3\ge4c-2b$\\
7 & $([8]\!\setminus\!\{8\};[8]\!\setminus\!\{5\};[8]\!\setminus\!\{5\})$ & $p_4\ge2c$\\
\bottomrule
\end{tabular}
\end{table}

\begin{proof}[Proof of Theorem~\ref{thm:phase}]
The entries in \eqref{eq:plow}, \eqref{eq:phigh}, and \eqref{eq:pface}
are nonnegative and weakly decreasing on their stated cones.  Exact Horn
feasibility is proved in Section~\ref{sec:replay}; Proposition~\ref{prop:horn}
therefore gives all three upper bounds.

For the first chamber, the weighted sum in Table~\ref{tab:lowdual} has left
side $\frac12\sum_{j=1}^7p_j$ and right side
\[
 \tfrac12(8c-6b)+\tfrac12(6c-6b)+\tfrac12(2c-2b)+2c
 =10c-7b.
\]
For the second chamber, Table~\ref{tab:highdual} gives
\[
 \tfrac12(8c-6b)+\tfrac12(4c-4b)+\tfrac12(2c)+2c
 =9c-5b.
\]
Thus the two global primals are optimal.

On $p_5=p_6=p_7=0$, half the sum of the four rows in
Table~\ref{tab:facedual} gives
\[
 \frac12\sum_{j=1}^4p_j\ge
 \tfrac12\{(8c-6b)+(6c-4b)+(4c-2b)+2c\}=10c-6b.
\]
The feasible vector \eqref{eq:pface} attains this bound.  Subtracting the two
global branches yields \eqref{eq:tent}.

For $0<b<c$, the gap is positive, so no rank-four factor lies on the global
optimal face.  Both global primals have exactly five positive entries there,
which proves $r_*=5$.  At $b=0$, \eqref{eq:plow} has rank four; at $b=c$,
\eqref{eq:phigh} has rank four.
\end{proof}

\section{Sharp tax and transition geometry}

\begin{corollary}[sharp relative rank-four tax]\label{cor:tax}
On the cone \eqref{eq:family},
\[
 1\le \frac{\kappa_8^{(\le4)}}{\kappa_8}\le\frac{14}{13}.
\]
The upper bound is attained exactly on the ray $b=c/2$.  The additive penalty
is at most $c/2$, with equality on the same ray.
\end{corollary}

\begin{proof}
After setting $t=b/c$, the ratio is
\[
 \frac{10-6t}{10-7t}\quad(0\le t\le1/2),\qquad
 \frac{10-6t}{9-5t}\quad(1/2\le t\le1).
\]
The first is strictly increasing and the second strictly decreasing.  Their
common value at $t=1/2$ is $14/13$.  Equation~\eqref{eq:tent} gives the
additive statement.
\end{proof}

The chamber wall is not where rank compatibility returns.  Both global
branches meet continuously at
\[
 p^L=p^H=(5c,3c,2c,2c,c,0,0),
\]
which still has rank five.  Instead, $b=c/2$ is where the identity of the
active Horn lower certificate switches and the rank penalty is largest.  Rank
four returns only on the two boundary rays, where the fifth singular square
vanishes.  The isolated spectrum $(11,3,3,3,-5,-5,-5,-5)$ corresponds to
$(b,c)=(3,5)$ and lies in the second chamber; Theorem~\ref{thm:phase} gives
$\kappa_8=30$, rank-four value $32$, and gap $2$ immediately.

\section{A two-valued separation in dimension nine}

The interior of the dimension-eight cone has three spectral levels.  The next
result shows that this multiplicity is not the mechanism behind the rank gap.

\begin{theorem}[two-valued rank gap]\label{thm:two-valued}
Let $G=G^*\in M_9(\mathbb C)$ have spectrum
\[
 (5,5,5,5,-4,-4,-4,-4,-4).
\]
Then $G$ has inertia $(4,5)$,
\[
 \kappa_9(G)=29,\qquad \kappa_9^{(\le5)}(G)=30,
\]
and the minimum rank on the unrestricted optimal face is six.
\end{theorem}

\begin{proof}
For the unrestricted problem and the rank-five face, respectively, use the
feasible squared-singular-value spectra
\[
 p=(16,12,10,10,8,2,0,0),\qquad
 p^{5}=(16,14,12,10,8,0,0,0).
\]
Their half-traces are $29$ and $30$.  Exact feasibility against all Horn,
order, and nonnegativity rows is part of the replay in
Section~\ref{sec:replay}.

For completeness, an unrestricted lower certificate uses weights
$1/2,1,1/2,1/2,1/2,1$ on the following five Horn triples and $p_8\ge0$:
\begin{align*}
 &(1;\{3\};\{1\};\{3\}),\quad
 (1;\{4\};\{1\};\{4\}),\\
 &(4;\{2,3,6,7\};\{1,2,6,7\};\{3,4,8,9\}),\\
 &(7;\{1,2,3,4,6,7,8\};\{1,2,3,4,6,7,8\};
       \{1,2,3,4,7,8,9\}),\\
 &(8;\{1,2,3,4,5,6,7,8\};\{1,2,3,4,6,7,8,9\};
       \{1,2,3,4,6,7,8,9\}).
\end{align*}
Their weighted sum is the full objective and has right side $29$.
On $p_6=p_7=p_8=0$, weights $1/2,1/2,1/2,1/2,3/2$ on
\begin{align*}
 &(1;\{4\};\{1\};\{4\}),\\
 &(4;\{3,4,6,7\};\{1,2,5,6\};\{3,4,8,9\}),\\
 &(6;\{2,3,4,6,7,8\};\{1,2,3,5,6,7\};\{2,3,4,7,8,9\}),\\
 &(7;\{1,2,3,4,6,7,8\};\{1,2,3,4,6,7,8\};
       \{1,2,3,4,7,8,9\}),\\
 &(8;\{1,2,3,4,5,6,7,8\};\{1,2,3,4,6,7,8,9\};
       \{1,2,3,4,6,7,8,9\})
\end{align*}
give the rank-five objective on that face and right side $30$.  Thus both
values are optimal.  Inertia gives the lower rank bound five.  The strict
rank-five gap rules out ranks at most five at the unrestricted optimum, while
the first displayed feasible vector has six positive entries.  Hence the
minimum optimal rank is six.
\end{proof}

\section{Exact replay and coverage}\label{sec:replay}

The dependency-free verifier regenerates Horn's recursive $T$-sets in
dimensions eight and nine.  In dimension eight, the counts by rank are
\[
 (36,462,2120,3516,2120,462,36),
\]
for $8{,}752$ Horn triples.  Six order rows and one nonnegativity row give
$8{,}759$ constraints.

Each right-hand side is affine in $t=b/c$, and each proposed primal vector is
affine on its chamber.  Hence every residual is affine.  Feasibility on
$[0,1/2]$ follows by checking $t=0$ and $t=1/2$; feasibility on $[1/2,1]$
follows from $t=1/2$ and $t=1$.  The rank-four branch is checked at $t=0$ and
$t=1$.  These endpoint checks are exact rational calculations, so they cover
the entire closed intervals rather than a grid.  The replay also verifies all
dual coefficient identities and their affine right-hand sides.

A second dimension-eight verifier independently enumerates nonzero
Littlewood--Richardson tableaux and checks the same parametric endpoints.  The
runner compares the family, chamber endpoints, Horn counts, primal branches,
dual values, rank statements and exact formulas.  No floating-point solver is
part of the proof artifact.  For Theorem~\ref{thm:two-valued}, the primary
replay regenerates
\[
 (45,792,5317,13704,13704,5317,792,45)
\]
Horn triples by rank: $39{,}716$ Horn rows and $39{,}724$ rows after order and
nonnegativity constraints.  It checks both primals and both dual identities in
exact rational arithmetic.

The replay certifies the finite Horn implications conditional on the classical
Horn theorem and Proposition~\ref{prop:horn}.  It does not constitute a formal
proof-assistant certification, a proof that dimension eight is minimal, a
classification of all dimension-eight spectra, priority adjudication, or peer
review.  The norm is the unnormalized Hilbert--Schmidt norm, matrices are
finite and complex, and the theorems concern only the displayed cone and
two-valued target.

\section{Conclusion}

Norm optimality and inertia-minimal rank separate on an entire explicit cone,
not merely at an exceptional target.  The global inverse cost changes Horn
facet at $b=c/2$, while the rank-four cost remains affine; their difference is
the tent function $\min\{b,c-b\}$.  This makes both the incompatibility and its
sharp size transparent.  The next structural problem is to determine whether
dimension eight is the first possible failure and to classify the larger Horn
chambers containing this cone.  The dimension-nine theorem separately shows
that repeated three-level structure is not necessary for the phenomenon.

\section*{Author and AI disclosure}

Lluis Eriksson directed the research, selected the claims, and accepts
responsibility for the manuscript.  OpenAI Codex assisted with conjecture
falsification, parametric Horn exploration, exact-certificate generation,
literature triage, adversarial proof checking, typesetting, and drafting.  No
external peer review or formal proof certification is claimed.

\AtBeginEnvironment{thebibliography}{\footnotesize\setlength{\itemsep}{0.2em}}
\bibliographystyle{plain}
\bibliography{references}

\end{document}
